\chapter*{术语表}

\begin{description}
\item[暴力法] 直接枚举解空间的方法，是理解题目和建立正确性的起点。
\item[不变量] 算法执行过程中始终成立的性质，用来证明优化解正确。
\item[滑动窗口] 用左右边界维护连续区间，并增量更新窗口状态的方法。
\item[前缀和] 把区间和查询转化为两个前缀差值的数据表示。
\item[哈希表] 通过哈希函数把 key 映射到桶，平均支持常数时间查询的数据结构。
\item[AC] Accepted，提交通过测试。AC 只是平台判定结果，不等于复杂度和边界已经讲清。
\item[LeetCode] 常用在线算法题库。本书引用题号时，题号是为了帮你定位完整题意和测试边界。
\item[LCP] 力扣杯竞赛题前缀，常见于力扣中文站。它表示题库来源，不表示算法类别。
\item[LCR] 力扣中国站重新编号的经典面试题集合前缀，很多题来自剑指 Offer 或专项题库。
\item[II/III/IV] 题名后的罗马数字，表示同一主题的后续变体。读题时要重新检查返回值、约束和构造要求。
\item[C++20] C++ 标准版本名。本书默认代码以 C++20 能力为边界，例如 \code{contains}、\code{string_view} 和部分 ranges 思想。
\item[C++23] 更新的 C++ 标准版本。平台未声明支持时，不应在刷题提交中默认使用。
\item[\code{std::}] C++ 标准库命名空间前缀。本书代码里的 \code{std::vector}、\code{std::string}、\code{std::queue}、\code{std::priority_queue} 分别对应动态数组、字符串、队列和堆等基础工具。
\item[\code{int64\_t}] 固定 64 位有符号整数，适合保存可能超过 \code{int} 的计数、区间和或乘积。
\item[\code{size\_t}] C++ 容器大小和下标常用类型。它是无符号类型，和负数混算时要小心。
\item[API] Application Programming Interface，接口。刷题里常指题目要求实现的方法或标准库调用合同。
\item[ID] Identifier，身份或编号。题目里的 id 通常用于唯一定位对象，不能和数组下标、值本身随意混用。
\item[IP] Internet Protocol，互联网协议地址。LeetCode 93 中的 IP 地址被简化成四段十进制数，每段 0 到 255 且无非法前导零。
\item[I/O] Input/Output，输入输出。算法题中通常指读入输出；工程题中还包括文件、网络和缓冲区边界。
\item[CPU] Central Processing Unit，中央处理器。算法复杂度最终要由 CPU 执行，但刷题阶段通常先用渐进复杂度建立主要判断。
\item[ASCII] 早期单字节字符编码集合。算法题若声明只含 ASCII，可以用固定大小计数数组。
\item[Unicode] 统一字符集合。工程文本处理不能简单把一个字节当一个用户可见字符。
\item[UTF-8] Unicode 的常见变长字节编码。LeetCode 393 这类题检查的是字节前缀模式和 continuation byte 合法性。
\item[STL] Standard Template Library，C++ 标准库容器、迭代器和算法的常见统称。
\item[INF] Infinity 的简写，表示“足够大”的哨兵值。使用时必须保证不会和真实答案或整数溢出混淆。
\item[EPSILON] 浮点比较容忍误差，用来处理二进制浮点舍入造成的极小偏差。
\item[DP] Dynamic Programming，动态规划。把重复子问题命名成状态，并保存状态答案。
\item[BFS] 广度优先搜索，适合无权最短路和层次扩展。
\item[DFS] 深度优先搜索，适合连通块、路径搜索和回溯。
\item[BST] Binary Search Tree，二叉搜索树。左小右大的有序树结构。
\item[AVL] 自平衡二叉搜索树，通过旋转控制高度差，保持对数级操作复杂度。
\item[DAG] Directed Acyclic Graph，有向无环图。适合拓扑排序、依赖分析和无环状态 DP。
\item[SCC] Strongly Connected Component，强连通分量。有向图中一组点两两互相可达。
\item[LCA] Lowest Common Ancestor，最近公共祖先。树上两个节点向根路径的最近交点。
\item[LRU] Least Recently Used，最近最少使用淘汰，按访问时间淘汰最旧元素。
\item[LFU] Least Frequently Used，最不经常使用淘汰，先按访问频次再按时间淘汰。
\item[LCS] Longest Common Subsequence，最长公共子序列。
\item[LIS] Longest Increasing Subsequence，最长递增子序列。
\item[KMP] Knuth-Morris-Pratt 字符串匹配算法，用失败函数避免主串回退。
\item[KM] Kuhn-Munkres 算法，常用于二分图最大权匹配。它处理的是带权匹配，不是普通无权最大匹配模板。
\item[DAWG] Directed Acyclic Word Graph，有向无环词图，是字符串集合或后缀关系的一类压缩表示。
\item[SPFA] Shortest Path Faster Algorithm，Bellman-Ford 的队列优化写法，复杂度不稳定。
\item[IDA*] Iterative Deepening A*，带启发式估价的迭代加深搜索。
\item[2-SAT] 二元布尔可满足性问题，常转成蕴含图并用强连通分量判断矛盾。
\item[mask/bitmask] 掩码或位掩码，用二进制位保存集合、访问状态或可用资源。
\item[URBG] Uniform Random Bit Generator，C++ 随机库中能产生随机比特序列的引擎概念。
\item[DNA] 脱氧核糖核酸。算法题里常抽象成只含 \code{A/C/G/T} 的字符串。
\item[IPO] Initial Public Offering，首次公开募股。对应题型常把项目选择建模成资本约束下的贪心和堆。
\item[FIFO] First In First Out，先进先出，典型结构是队列。
\item[LIFO] Last In First Out，后进先出，典型结构是栈。
\item[RCU] Read-Copy-Update，读路径尽量无锁、更新路径复制并延迟回收的同步思想。
\item[SLAB/SLUB] Linux 内核小对象分配器族，用固定大小对象缓存降低分配成本。SLUB 是现代 Linux 常见实现。
\item[VMA] Virtual Memory Area，Linux 进程地址空间中的一段连续虚拟内存区域。
\item[XArray] Linux 内核中的稀疏整数索引结构，用整数 id 查对象。
\item[状态常量] 代码中给题目状态取的名字，例如 \code{FRESH}、\code{ROTTEN}、\code{LAND}、\code{WATER}。它们不是算法名，而是让分支条件更贴近题意。
\item[哨兵常量] 人为设置的特殊初值或边界值，例如 \code{NEG\_INF}、\code{INF}、\code{MAX\_VALUE}。它们必须和题目真实值域分开。
\item[字符集容量常量] 计数字符时使用的数组容量，例如 \code{ALPHABET\_SIZE} 或 \code{ASCII\_RANGE}。容量来自题目字符集，而不是固定模板。
\item[样例字符串] 题目示例中用于手推的具体字符串，例如 \code{ADOBECODEBANC}。它不是缩写，作用是让你观察窗口、匹配或递归状态如何变化。
\item[命名常量] 用名字替代魔法数字的代码常量，例如 \code{DEFAULT_CAPACITY}、\code{SEGMENT_COUNT}。理解它时要回到当前题目的约束或状态设计。
\item[CSR] Compressed Sparse Row，压缩稀疏行格式，常用于图和稀疏矩阵。
\item[NUMA] Non-Uniform Memory Access，非统一内存访问，不同内存节点访问成本不同。
\item[拓扑排序] 对有向无环图按依赖顺序排列节点的方法。
\end{description}
